\documentclass[10pt,a4paper]{amsart}
\usepackage{amsmath,amssymb,amsthm,graphicx}
% Retained preamble for the Stacks Project chapter {CH} (shared project preamble).
% The chapter file pinned-original-source.tex only carries \input{{preamble}} on its first line,
% so no preamble text lives inside the chapter itself.  The Stacks Project uses one single shared
% preamble for every chapter: the snapshot stages it three times (sources/stacks-chow/preamble.tex,
% sources/stacks-etale-cohomology/preamble.tex and focused-seeds/AGP/sources/stacks-snapshot-20261011/
% preamble.tex) and those copies are identical in size.  The block below is copied mechanically,
% line for line, from sources/stacks-chow/preamble.tex lines 14-18 and 154-217: the Xy-pic setup and
% the native theorem declarations together with the mathematical macros.  Nothing was added, removed
% or redefined; the wrapper class and the amsmath/amssymb/amsthm/graphicx base packages are supplied
% by the assembler.
% For commutative diagrams we use Xy-pic
% Chapter: stacks-spaces-more-morphisms.tex -- More on Morphisms of Spaces.
\usepackage[all]{xy}

% We use 2cell for 2-commutative diagrams.
\xyoption{2cell}

% Theorem environments.
%
\theoremstyle{plain}
\newtheorem{theorem}[subsection]{Theorem}
\newtheorem{proposition}[subsection]{Proposition}
\newtheorem{lemma}[subsection]{Lemma}

\theoremstyle{definition}
\newtheorem{definition}[subsection]{Definition}
\newtheorem{example}[subsection]{Example}
\newtheorem{exercise}[subsection]{Exercise}
\newtheorem{situation}[subsection]{Situation}

\theoremstyle{remark}
\newtheorem{remark}[subsection]{Remark}
\newtheorem{remarks}[subsection]{Remarks}

\numberwithin{equation}{subsection}

% Macros
%
\def\lim{\mathop{\mathrm{lim}}\nolimits}
\def\colim{\mathop{\mathrm{colim}}\nolimits}
\def\Spec{\mathop{\mathrm{Spec}}}
\def\Hom{\mathop{\mathrm{Hom}}\nolimits}
\def\Ext{\mathop{\mathrm{Ext}}\nolimits}
\def\SheafHom{\mathop{\mathcal{H}\!\mathit{om}}\nolimits}
\def\SheafExt{\mathop{\mathcal{E}\!\mathit{xt}}\nolimits}
\def\Sch{\mathit{Sch}}
\def\Mor{\mathop{\mathrm{Mor}}\nolimits}
\def\Ob{\mathop{\mathrm{Ob}}\nolimits}
\def\Sh{\mathop{\mathit{Sh}}\nolimits}
\def\NL{\mathop{N\!L}\nolimits}
\def\CH{\mathop{\mathrm{CH}}\nolimits}
\def\proetale{{pro\text{-}\acute{e}tale}}
\def\etale{{\acute{e}tale}}
\def\QCoh{\mathit{QCoh}}
\def\Ker{\mathop{\mathrm{Ker}}}
\def\Im{\mathop{\mathrm{Im}}}
\def\Coker{\mathop{\mathrm{Coker}}}
\def\Coim{\mathop{\mathrm{Coim}}}

% Boxtimes
%
\DeclareMathSymbol{\boxtimes}{\mathbin}{AMSa}{"02}

%
% Macros for moduli stacks/spaces
%
\def\QCohstack{\mathcal{QC}\!\mathit{oh}}
\def\Cohstack{\mathcal{C}\!\mathit{oh}}
\def\Spacesstack{\mathcal{S}\!\mathit{paces}}
\def\Quotfunctor{\mathrm{Quot}}
\def\Hilbfunctor{\mathrm{Hilb}}
\def\Curvesstack{\mathcal{C}\!\mathit{urves}}
\def\Polarizedstack{\mathcal{P}\!\mathit{olarized}}
\def\Complexesstack{\mathcal{C}\!\mathit{omplexes}}
% \Pic is the operator that assigns to X its picard group, usage \Pic(X)
% \Picardstack_{X/B} denotes the Picard stack of X over B
% \Picardfunctor_{X/B} denotes the Picard functor of X over B
\def\Pic{\mathop{\mathrm{Pic}}\nolimits}
\def\Picardstack{\mathcal{P}\!\mathit{ic}}
\def\Picardfunctor{\mathrm{Pic}}
\def\Deformationcategory{\mathcal{D}\!\mathit{ef}}
\usepackage{cleveref}
\crefname{originallemma}{original lemma}{original lemmas}
\Crefname{originallemma}{Original lemma}{Original lemmas}
\makeatletter
\newlabel{descent-definition-descent-datum-modules}{{023G}{source}{}{sourcepointer-1}{}}
\newlabel{descent-definition-descent-datum-modules@cref}{{[originallemma][0][]023G}{[1][0][]source}{}{}{}}
\makeatother
\crefname{originallemma}{original lemma}{original lemmas}
\Crefname{originallemma}{Original lemma}{Original lemmas}
\makeatletter
\newlabel{descent-subsection-descent-modules-morphisms}{{08WW}{source}{}{sourcepointer-2}{}}
\newlabel{descent-subsection-descent-modules-morphisms@cref}{{[originallemma][0][]08WW}{[1][0][]source}{}{}{}}
\makeatother
\crefname{originallemma}{original lemma}{original lemmas}
\Crefname{originallemma}{Original lemma}{Original lemmas}
\makeatletter
\newlabel{descent-theorem-descent}{{08XA}{source}{}{sourcepointer-3}{}}
\newlabel{descent-theorem-descent@cref}{{[originallemma][0][]08XA}{[1][0][]source}{}{}{}}
\makeatother
\crefname{originallemma}{original lemma}{original lemmas}
\Crefname{originallemma}{Original lemma}{Original lemmas}
\makeatletter
\newlabel{spaces-descent-lemma-descending-property-isomorphism}{{041Y}{source}{}{sourcepointer-4}{}}
\newlabel{spaces-descent-lemma-descending-property-isomorphism@cref}{{[originallemma][0][]041Y}{[1][0][]source}{}{}{}}
\makeatother
\crefname{originallemma}{original lemma}{original lemmas}
\Crefname{originallemma}{Original lemma}{Original lemmas}
\makeatletter
\newlabel{spaces-properties-lemma-quasi-compact-affine-cover}{{03H6}{source}{}{sourcepointer-5}{}}
\newlabel{spaces-properties-lemma-quasi-compact-affine-cover@cref}{{[originallemma][0][]03H6}{[1][0][]source}{}{}{}}
\makeatother
\crefname{originallemma}{original lemma}{original lemmas}
\Crefname{originallemma}{Original lemma}{Original lemmas}
\makeatletter
\newlabel{spaces-properties-proposition-quasi-coherent}{{03M3}{source}{}{sourcepointer-6}{}}
\newlabel{spaces-properties-proposition-quasi-coherent@cref}{{[originallemma][0][]03M3}{[1][0][]source}{}{}{}}
\makeatother
\title{A quasi-compact monomorphism of algebraic spaces with injective structure sheaf pullbacks is an isomorphism}
\author{The Stacks Project Authors}
\date{stacks-project:spaces-more-morphisms.tex:local-snapshot-20261011; GNU FDL 1.2 or later}
\begin{document}
\maketitle
\noindent\textit{Editorial note.} All delivered mathematical text is the author's own, copied from the pinned Stacks Project chapter file spaces-more-morphisms.tex; no statement or proof step was rewritten, reordered or removed.\par\medskip
\noindent\textit{Editorial note.} Preamble provenance, as required for this source: the first line of the pinned chapter file is \textbackslash{}input\{preamble\}, and the native preamble was not staged with the source, so no preamble text lives inside the chapter. The retained preamble used here is the project's single Stacks-wide preamble (the same 217-line / 8116-byte file that the already published chow and etale-cohomology cards use), copied mechanically line for line from sources/stacks-chow/preamble.tex lines 14-18 (Xy-pic setup) and 154-217 (theorem declarations and mathematical macros). Not one character was changed, no package was added, no macro was defined, and the excerpt serves only as preamble, never as a span coordinate.\par\medskip
\noindent\textit{Editorial note.} The same-chapter pointers are located at their \textbackslash{}label lines inside the pinned chapter file; the cross-chapter pointers are located in the project's own tag list staged in the snapshot (sources/stacks-chow/tags.txt), at the verbatim lines carrying their canonical tags.\par\medskip
\noindent\textit{Allowed standard prerequisites (editorial).}\begin{itemize}
\item Existence of an affine etale cover of a quasi-compact algebraic space. Properties of Spaces, Lemma labelled spaces-properties-lemma-quasi-compact-affine-cover; invoked at lines 331-333. Supplies U = Spec(B) with a surjective etale morphism U -> X.
\item Quasi-coherent modules on X are descent data for A -> B. Properties of Spaces, Proposition labelled spaces-properties-proposition-quasi-coherent together with Descent, Definition labelled descent-definition-descent-datum-modules; invoked at lines 335-340. Identifies the category of quasi-coherent O\_X-modules with DD\_\{B/A\}.
\item Descent for modules along a universally injective ring map. Descent, Theorem labelled descent-theorem-descent; invoked at line 349. Reduces DD\_\{B/A\} to the category of A-modules once A -> B is seen to be universally injective.
\item fpqc descent of the property of being an isomorphism. Descent on Spaces, Lemma labelled spaces-descent-lemma-descending-property-isomorphism; invoked at lines 358-359. Turns 'becomes an isomorphism after base change by X -> Spec(A)' into 'is an isomorphism'.
\end{itemize}\medskip
\ifdefined\hypertarget\hypertarget{sourcepointer-1}{}\fi
\noindent\textit{Source reference.} Original lemma 023G: sources/stacks-chow/tags.txt, source lines 2741--2741. Navigation pointer for \textbackslash{}ref\{descent-definition-descent-datum-modules\} used inside the counted span (line 339); the argument that consumes it is retained in the span.\par\medskip
\ifdefined\hypertarget\hypertarget{sourcepointer-2}{}\fi
\noindent\textit{Source reference.} Original lemma 08WW: sources/stacks-chow/tags.txt, source lines 11284--11284. Navigation pointer for \textbackslash{}ref\{descent-subsection-descent-modules-morphisms\} used inside the counted span (line 340); the argument that consumes it is retained in the span.\par\medskip
\ifdefined\hypertarget\hypertarget{sourcepointer-3}{}\fi
\noindent\textit{Source reference.} Original lemma 08XA: sources/stacks-chow/tags.txt, source lines 11298--11298. Navigation pointer for \textbackslash{}ref\{descent-theorem-descent\} used inside the counted span (line 349); the argument that consumes it is retained in the span.\par\medskip
\ifdefined\hypertarget\hypertarget{sourcepointer-4}{}\fi
\noindent\textit{Source reference.} Original lemma 041Y: sources/stacks-chow/tags.txt, source lines 5309--5309. Navigation pointer for \textbackslash{}ref\{spaces-descent-lemma-descending-property-isomorphism\} used inside the counted span (line 359); the argument that consumes it is retained in the span.\par\medskip
\ifdefined\hypertarget\hypertarget{sourcepointer-5}{}\fi
\noindent\textit{Source reference.} Original lemma 03H6: sources/stacks-chow/tags.txt, source lines 4555--4555. Navigation pointer for \textbackslash{}ref\{spaces-properties-lemma-quasi-compact-affine-cover\} used inside the counted span (line 333); the argument that consumes it is retained in the span.\par\medskip
\ifdefined\hypertarget\hypertarget{sourcepointer-6}{}\fi
\noindent\textit{Source reference.} Original lemma 03M3: sources/stacks-chow/tags.txt, source lines 4735--4735. Navigation pointer for \textbackslash{}ref\{spaces-properties-proposition-quasi-coherent\} used inside the counted span (line 338); the argument that consumes it is retained in the span.\par\medskip
\begin{lemma}
\label{lemma-ui-case}
Let $S$ be a scheme. Let $f : X \to Y$ be a quasi-compact monomorphism
of algebraic spaces such that for every $T \to Y$ the map
$$
\mathcal{O}_T \to f_{T,*}\mathcal{O}_{X \times_Y T}
$$
is injective. Then $f$ is an isomorphism (and hence representable by schemes).
\end{lemma}

\begin{proof}
The question is \'etale local on $Y$, hence we may assume $Y = \Spec(A)$
is affine. Then $X$ is quasi-compact and we may choose an affine scheme
$U = \Spec(B)$ and a surjective \'etale morphism $U \to X$
(Properties of Spaces, Lemma
\ref{spaces-properties-lemma-quasi-compact-affine-cover}).
Note that $U \times_X U = \Spec(B \otimes_A B)$. Hence the category of
quasi-coherent $\mathcal{O}_X$-modules is equivalent to the
category $DD_{B/A}$ of descent data on modules for $A \to B$.
See Properties of Spaces, Proposition
\ref{spaces-properties-proposition-quasi-coherent},
Descent, Definition \ref{descent-definition-descent-datum-modules}, and
Descent, Subsection \ref{descent-subsection-descent-modules-morphisms}.
On the other hand,
$$
A \to B
$$
is a universally injective ring map. Namely, given an
$A$-module $M$ we see that $A \oplus M \to B \otimes_A (A \oplus M)$
is injective by the assumption of the lemma. Hence
$DD_{B/A}$ is equivalent to the category of $A$-modules by
Descent, Theorem \ref{descent-theorem-descent}. Thus pullback along
$f : X \to \Spec(A)$ determines an equivalence of categories of
quasi-coherent modules. In particular $f^*$ is exact on
quasi-coherent modules and we see that $f$ is flat
(small detail omitted). Moreover, it is clear that $f$ is surjective
(for example because $\Spec(B) \to \Spec(A)$ is surjective).
Hence we see that $\{X \to \Spec(A)\}$ is an fpqc cover.
Then $X \to \Spec(A)$ is a morphism which becomes an isomorphism
after base change by $X \to \Spec(A)$. Hence it is an isomorphism by
fpqc descent, see Descent on Spaces, Lemma
\ref{spaces-descent-lemma-descending-property-isomorphism}.
\end{proof}

\end{document}
